\section{Studi Kasus: Basis Pengetahuan Silsilah Keluarga}

Subsection ini menyajikan studi kasus lengkap: merepresentasikan silsilah keluarga dengan fakta dan aturan Prolog, lalu menjalankan berbagai query \cite{learn_prolog,swi_manual}. Tujuan studi kasus ini adalah mengintegrasikan materi Sub-CPMK 4.3: menulis fakta dan aturan serta menjalankan query dasar.

\subsection{Requirements}

Basis pengetahuan silsilah harus memuat:
\begin{itemize}[leftmargin=*, nosep]
  \item \textbf{Fakta} \texttt{orangtua(OrangTua, Anak)}: relasi siapa orang tua dari siapa (minimal 6 fakta).
  \item \textbf{Fakta} \texttt{laki(X)} dan \texttt{perempuan(X)}: jenis kelamin (opsional tetapi berguna untuk ayah/ibu).
  \item \textbf{Aturan} \texttt{ayah(X, Y)}: X ayah dari Y (berdasarkan \texttt{orangtua} dan \texttt{laki}).
  \item \textbf{Aturan} \texttt{ibu(X, Y)}: X ibu dari Y (berdasarkan \texttt{orangtua} dan \texttt{perempuan}).
  \item \textbf{Aturan} \texttt{kakek(X, Z)}: X kakek dari Z (ayah dari orang tua Z).
  \item \textbf{Aturan} \texttt{nenek(X, Z)}: X nenek dari Z (ibu dari orang tua Z).
\end{itemize}

\subsection{Program Prolog Lengkap}

Berikut implementasi lengkap yang memenuhi requirements di atas \cite{learn_prolog,swi_prolog}.

\begin{lstlisting}[style=prologstyle, caption={Program silsilah keluarga: fakta dan aturan}]
% Fakta: orangtua(OrangTua, Anak)
orangtua(hasan, ahmad).
orangtua(hasan, rina).
orangtua(minah, ahmad).
orangtua(minah, rina).
orangtua(ahmad, ali).
orangtua(ahmad, sari).
orangtua(rina, budi).

% Fakta: jenis kelamin
laki(hasan).
laki(ahmad).
laki(ali).
laki(budi).
perempuan(minah).
perempuan(rina).
perempuan(sari).

% Aturan: ayah, ibu, kakek, nenek
ayah(X, Y) :- orangtua(X, Y), laki(X).
ibu(X, Y) :- orangtua(X, Y), perempuan(X).
kakek(X, Z) :- ayah(X, Y), orangtua(Y, Z).
nenek(X, Z) :- ibu(X, Y), orangtua(Y, Z).
\end{lstlisting}

\subsection{Contoh Query dan Jawaban}

Setelah program dimuat di SWI-Prolog, query berikut dapat dijalankan. Tabel~\ref{tab:studi-kasus-query} merangkum beberapa query dan jawaban yang diharapkan \cite{learn_prolog,swi_manual}.

\begin{table}[htbp]
\centering
\small
\begin{tabular}{|>{\raggedright\arraybackslash}p{4.2cm}|>{\raggedright\arraybackslash}p{6cm}|}
\hline
\textbf{Query} & \textbf{Jawaban (contoh)} \\
\hline
\texttt{?- kakek(X, Y).} & \texttt{X = hasan, Y = ali} ; \texttt{X = hasan, Y = sari} ; \texttt{X = hasan, Y = budi} \\
\hline
\texttt{?- nenek(X, ali).} & \texttt{X = minah} \\
\hline
\texttt{?- orangtua(ahmad, X).} & \texttt{X = ali} ; \texttt{X = sari} \\
\hline
\texttt{?- ayah(hasan, X).} & \texttt{X = ahmad} ; \texttt{X = rina} \\
\hline
\texttt{?- kakek(hasan, budi).} & \texttt{true} \\
\hline
\end{tabular}
\caption{Contoh query studi kasus silsilah dan jawaban \cite{learn_prolog}}
\label{tab:studi-kasus-query}
\end{table}

Contoh interaksi di konsol Prolog:

\begin{lstlisting}[style=prologstyle, caption={Contoh menjalankan query silsilah}]
?- kakek(X, Y).
X = hasan,
Y = ali ;
X = hasan,
Y = sari ;
X = hasan,
Y = budi.

?- nenek(X, ali).
X = minah.
\end{lstlisting}

\subsection{Kaitan dengan Logika Predikat}

Program silsilah di atas dapat dipandang sebagai teori dalam logika predikat orde pertama \cite{olp_fol,berkeley_fol}. Setiap fakta adalah atom (misalnya \texttt{orangtua(hasan, ahmad)}), dan setiap aturan bersesuaian dengan klausa Horn: implikasi dengan satu literal positif di kepala dan konjungsi literal di tubuh. Misalnya aturan \texttt{kakek(X,Z) :- ayah(X,Y), orangtua(Y,Z).} dalam notasi logika: $\forall X\,\forall Y\,\forall Z\;\big(\big(\mathrm{ayah}(X,Y) \wedge \mathrm{orangtua}(Y,Z)\big) \rightarrow \mathrm{kakek}(X,Z)\big)$. Mesin Prolog menggunakan unifikasi dan SLD resolution untuk menjawab query; dengan demikian, menjalankan query sama dengan meminta sistem membuktikan apakah goal tersebut mengikuti dari teori (fakta dan aturan) \cite{learn_prolog,lpn_unification}.

\begin{konsep}
Studi kasus silsilah: fakta \texttt{orangtua/2}, \texttt{laki/1}, \texttt{perempuan/1}; aturan \texttt{ayah/2}, \texttt{ibu/2}, \texttt{kakek/2}, \texttt{nenek/2}. Query memanfaatkan variabel untuk menanyakan relasi. Program bersesuaian dengan klausa Horn dalam logika predikat.
\end{konsep}
